% Recipe: xelatex -> bibtex -> xelatex*2

\documentclass{mySimpleArticleCN}

\title{\textbf{\LaTeX{}排版示例}\thanks{本文档受到乌有先生之乌有基金的资助}}
\author{You Know Who, \quad You Cannot-Know Who\thanks{小组长/指导教师/通讯作者：wuyou@nosuchanemail.no}}
\date{\today}

\begin{document}
    \maketitle          % 生成封面
    \tableofcontents    % 目录
    \newpage
    
    \section{图表示例}

    \subsection{插入图片}

    如图~\ref{fig:starryNight}~所示。图片存放在 \verb|figures| 目录。

    \begin{figure}[htbp]
        \centering
        \includegraphics[width=0.6\linewidth]{starryNight.jpg}
        \caption{A figure of starry night by Vincent}
        \label{fig:starryNight}
    \end{figure}

    \subsection{插入表格}

    如表~\ref{tab:Univ-rank-2024}~所示。

    \begin{table}[htbp]
        \centering
        \caption{Chinese Universities Ranked in the Top 50 in QS 2024}
        \label{tab:Univ-rank-2024}
        \begin{tabular}{@{}ccc@{}}
        \toprule
        Rank & University Name                               & Location         \\ \midrule
        17   & Peking University                             & China (Mainland) \\
        25   & \textbf{Tsinghua University}                  & China (Mainland) \\
        44   & Zhejiang University                           & China (Mainland) \\
        50   & Fudan University                              & China (Mainland) \\
        \bottomrule
        \end{tabular}
    \end{table}


    \section{数学公式示例}

    在量子力学中，微观粒子在势场 \(V(\mathbf{r})\) 中的运动可以描述为
    \[
        \mathrm{i} \hbar \frac{\partial \Psi}{\partial t} = 
        - \frac{\hbar^2}{2m} \nabla^2 \Psi + V(\mathbf{r}) \Psi
    \]

    Maxwell电磁方程组如\eqref{eq:maxwell-equations}所示。
    \begin{equation}
        \label{eq:maxwell-equations}
        \left\{
        \begin{aligned}
            \nabla \cdot \mathbf{D} &= \rho_0 \\
            \nabla \cdot \mathbf{B} &= 0 \\
            \nabla \times \mathbf{E} &= -\frac{\partial \mathbf{B}}{\partial t} \\
            \nabla \times \mathbf{H} &= \mathbf{j}_0 + \frac{\partial \mathbf{D}}{\partial t}
        \end{aligned}
        \right.
    \end{equation}

    Maxwell电磁方程组如\eqref{eq:Maxwell-Eqns-1}--\eqref{eq:Maxwell-Eqns-4}所示。
    \begin{align}
        \label{eq:Maxwell-Eqns-1}
        \nabla \cdot \mathbf{D} &= \rho_0 \\
        \label{eq:Maxwell-Eqns-2}
        \nabla \cdot \mathbf{B} &= 0 \\
        \label{eq:Maxwell-Eqns-3}
        \nabla \times \mathbf{E} &= -\frac{\partial \mathbf{B}}{\partial t} \\
        \label{eq:Maxwell-Eqns-4}
        \nabla \times \mathbf{H} &= \mathbf{j}_0 + \frac{\partial \mathbf{D}}{\partial t}
    \end{align}

    可以用 \verb|subequations| 环境使不同方程拥有公共编号和各自的子编号。例如：
    \begin{subequations}
        \begin{align}
            \label{eq:Maxwell-Eqns-01}
            \nabla \cdot \mathbf{D} &= \rho_0 \\
            \label{eq:Maxwell-Eqns-02}
            \nabla \cdot \mathbf{B} &= 0 \\
            \label{eq:Maxwell-Eqns-03}
            \nabla \times \mathbf{E} &= -\frac{\partial \mathbf{B}}{\partial t} \\
            \label{eq:Maxwell-Eqns-04}
            \nabla \times \mathbf{H} &= \mathbf{j}_0 + \frac{\partial \mathbf{D}}{\partial t}
        \end{align}
    \end{subequations}

    可使用 \verb|cases| 环境产生分段函数，例如：
    \begin{equation}
        f(x,y)=
        \begin{cases}
        \dfrac{xy}{x^2+y^2}, & x^2+y^2 \neq 0 \\
        0, & x=y=0
        \end{cases}
    \end{equation}


    \section{参考文献引用示例}
        \subsection{引言}

        \lipsum[1]\cite{jack1971,lieberman1989}

        \lipsum[3]\cite{sander1969}

        \subsection{方法}

        \lipsum[4]\cite{Powis2024,Sun2023a}

        \subsection{结果与讨论}

        \lipsum[5-6]\citep{Lopez2013, Falgout2002a}

        \lipsum[7]\citep{Donko2021}

        \subsection{结论}

        \lipsum[8]

\newpage

% 列出参考文献
\bibliography{reference/refs}   % BibTeX 数据库文件(.bib)

\end{document}
